% संपूर्ण मराठी ग्रंथातील प्रकरण 66: निवड
% हा संपादनयोग्य स्रोतखंड आहे. संकलनासाठी संपूर्ण स्रोतसंग्रहातील
% अक्षररूपे, पूर्वभाग, चिन्हव्याख्या आणि आकृती-स्रोत आवश्यक आहेत.
% मूळ: Open Logic Project; मजकूर आणि रूपांतर CC BY 4.0.
% यंत्रानुवाद, दुरुस्ती आणि तपासणी: OpenAI Codex —
% GPT-5.6 Sol आणि GPT-6 Sol, दोन्ही Ultra effort.
% दुरुस्ती-पुनरावलोकन: GPT-6.1 Sol, Ultra effort.
\chapter{निवड}\label{sth:choice:chap}





\section{प्रस्तावना}\label{sth:choice:intro:sec}

\ref{sth:cardinals:chap}--\ref{sth:card-arithmetic:chap} मध्ये संचांकांना क्रमसूचक संख्या मानून संचांकांची उपपत्ती विकसित केली. ही पद्धत सुक्रमणाच्या स्वयंसिद्धकावर अवलंबून आहे. सुक्रमण हे दुसऱ्या एका तत्त्वाशी---निवडीच्या स्वयंसिद्धकाशी---सममूल्य ठरते; आणि त्या तत्त्वाच्या स्वीकारार्हतेविषयी गंभीर तात्त्विक चर्चा झाली आहे. या प्रकरणातील आपले प्रश्न असे आहेत: हे स्वयंसिद्धक कसे वापरले जाते आणि त्याचे समर्थन करता येते का?







\section{टार्स्की--स्कॉटची युक्ती}\label{sth:choice:tarskiscott:sec}

\ref{sth:cardinals:cardsasords:defcardinalasordinal} मध्ये संचांकांना क्रमसूचक संख्या म्हणून परिभाषित केले. यासाठी सुक्रमणाचे स्वयंसिद्धक गृहीत धरले होते. हीच प्रचलित प्रमाणपद्धत आहे, एवढ्याच कारणासाठी आपण तसे केले.

म्हणून सुक्रमणाशी संबंधित तात्त्विक प्रश्नांची चर्चा करण्याआधी, प्रचलित प्रमाणपद्धत सोडून सुक्रमण गृहीत \emph{न धरता} संचांकांची उपपत्ती विकसित करणे आपल्याला \emph{शक्य} आहे, हे स्पष्ट असणे महत्त्वाचे आहे. \ref{sfr:siz::chap} मध्ये आलेल्या $\cardeq{A}{B}$, $\cardle{A}{B}$ आणि $\cardless{A}{B}$ या व्याख्या तरीही वापरता येतील. फक्त \emph{संचांकाची} नवी संकल्पना लागेल.

$A$ चा संचांक पुढीलप्रमाणे परिभाषित करावा, अशी साधी कल्पना सुचू शकते:
\[
	\Setabs{x}{\cardeq{A}{x}}.
\]
\ref{sth:cardinals:hp:sec} च्या अगदी शेवटी मांडलेल्या फ्रेगेच्या $\# x Fx$ या व्याख्येशी याची तुलना करता येईल. तिथे सूचित केलेल्या कारणांमुळे ही सर्वसाधारण व्याख्या अपयशी ठरते. कोणताही एकघटकी संच $\{\emptyset\}$ शी तुल्यबल असतो. पण श्रेणीच्या प्रत्येक उत्तरवर्ती टप्प्यावर नवे एकघटकी संच तयार होतात (आधीच्या टप्प्याचा एकघटकी संचच पाहा). त्यामुळे अरिक्त आधारसंचाच्या बाबतीत $\Setabs{x}{\cardeq{A}{x}}$ अस्तित्वात नसतो, कारण त्याला प्रतांक असू शकत नाही. रिक्त आधारसंचाचे प्रकरण या आक्षेपाच्या कक्षेत येत नाही.

ही अडचण टाळण्यासाठी टार्स्की आणि स्कॉट यांची युक्ती वापरू:\footnote{आठवण: स्पष्टपणे वेगळे सांगितले नसेल, तर सर्व सूत्रांत प्राचल असू शकतात.}

\begin{defn}[टार्स्की--स्कॉट]
कोणत्याही सूत्र $\phi(x)$ साठी, $\phi(x)$ पूर्ण करणाऱ्या आणि शक्य तितक्या लघुतम प्रतांकाच्या सर्व $x$ चा संच $[ x : \phi(x)] $ असू द्या ($\phi$ पूर्ण करणारी कोणतीही वस्तू नसल्यास $\emptyset$ घ्या).
\end{defn}

ही व्याख्या योग्य आहे हे तपासले पाहिजे. $\ZF$ मध्ये काम करताना \ref{sth:spine:foundation:zfentailsregularity} नुसार प्रत्येक $x$ साठी $\setrank{x}$ अस्तित्वात असतो. आता $\phi$ पूर्ण करणाऱ्या काही वस्तू असतील, तर $(\exists x\subseteq V_\alpha)\phi(x)$ पूर्ण करणारा लघुतम प्रतांक $\alpha$ घेता येतो; त्याच्या लघुतमतेचा अर्थ $(\forall \beta
\in \alpha)(\forall x \subseteq V_\beta)\lnot \phi(x)$ असा आहे. त्यानंतर विभक्तीकरणाने $[x : \phi(x)]$ हा संच $\Setabs{x \in
V_{\alpha+1}}{\phi(x)}$ म्हणून परिभाषित करता येतो.

टार्स्की--स्कॉटच्या युक्तीचे समर्थन केले आहे; आता ती वापरून संचांकाची संकल्पना परिभाषित करता येईल:

\begin{defn}
$A$ चा \textsc{ts}-संचांक म्हणजे $\text{tsc}(A) = [x :
\cardeq{A}{x}]$.
\end{defn}

\textsc{ts}-संचांकाच्या व्याख्येत सुक्रमण वापरलेले नाही. पण ते स्वयंसिद्धक नसतानाही \emph{\textsc{ts}-संचांकांचे} वर्तन \ref{sth:cardinals:cardsasords:defcardinalasordinal} मध्ये परिभाषित केलेल्या \emph{संचांकांसारखेच} असल्याचे दाखवता येते. उदाहरणार्थ, \ref{sth:cardinals:cardsasords:lem:CardinalsBehaveRight} आणि \ref{sth:card-arithmetic:opps:lem:SizePowerset2Exp} हे निष्कर्ष \textsc{ts}-संचांकांच्या भाषेत पुन्हा मांडल्यास, सुक्रमण गृहीत न धरता त्यांचे पुरावे $\ZF$ मध्ये व्यवस्थित चालतात.

याच संदर्भात, टार्स्की--स्कॉटची युक्ती वापरून क्रमसूचक संख्यांची उपपत्तीही विकसित करता येते. $\tuple{A, <}$ हे सुक्रमण असताना $\text{tso}(A, <) = [\tuple{X,
R} : \ordeq{\tuple{A, <}}{\tuple{X, R}}]$ घ्या. संचांक आणि क्रमसूचक संख्या यांच्या या मांडणीविषयी अधिक माहितीसाठी पॉटर (2004, प्रकरणे~9--12) पाहा.







\section{तुलनाक्षमता आणि हार्टोग्सचे पूर्वप्रमेय}\label{sth:choice:hartogs:sec}

ही झाली अनुकूल बाजू. आता प्रतिकूल बाजू पाहू. निवडीचे स्वयंसिद्धक नसल्यास गोष्टी \emph{गुंतागुंतीच्या} होतात. असे का, हे पाहण्यासाठी (हार्टोग्स 1915) यांचा हा सुंदर निष्कर्ष पाहा:

\begin{lem}[\emph{$\ZF$ मध्ये}]\label{sth:choice:hartogs:HartogsLemma}
कोणत्याही संच $A$ साठी $\cardnless{\alpha}{A}$ असलेली क्रमसूचक संख्या $\alpha$ अस्तित्वात आहे.
\end{lem}

\begin{proof}
$B \subseteq A$ आणि $R \subseteq B^2$ असल्यास \ref{sth:ord-arithmetic:using-addition:rankcomputation} नुसार $\tuple{B, R} \subseteq
V_{\setrank{A}+4}$. म्हणून विभक्तीकरण वापरून पुढील संच पाहू:
\[
	C = \Setabs{\tuple{B, R} \in V_{\setrank{A}+5}}{B\subseteq A
	\text{ आणि $\tuple{B, R}$ हे सुक्रमण आहे}}
\]
प्रतिस्थापन आणि \ref{sth:ordinals:ordtype:thmOrdinalRepresentation} वापरून पुढील संच तयार करा:
\[
	\alpha = \Setabs{\ordtype{B, R}}{\tuple{B, R} \in C}.
\]
\ref{sth:ordinals:basic:corordtransitiveord} नुसार $\alpha$ ही क्रमसूचक संख्या आहे, कारण तो क्रमसूचक संख्यांचा संक्रमक संच आहे. कारण $\gamma \in \beta \in \alpha$ असेल, तर एखाद्या संचासाठी $\beta = \ordtype{B, R}$, जेथे $B \subseteq A$, आणि मग \ref{sth:ordinals:iso:wellordinitialsegment} नुसार एखाद्या $b
\in B$ साठी $\gamma = \ordtype{B_b, R_b}$. त्यामुळे $\gamma \in \alpha$.


विरोधसिद्धीसाठी एकास-एक फलन $f \colon \alpha \to A$ अस्तित्वात आहे असे समजा. मग पुढील व्याख्या करा:
\begin{align*}
	B &= \ran{f}\\
	R &= \Setabs{\tuple{f(\xi), f(\eta)} \in A \times A}{\xi \in \eta \land \xi \in \alpha \land \eta \in \alpha}.
\end{align*}
येथे अंतर्गत निर्देशांक फलनाच्या प्रांतात घेतले आहेत. स्पष्टच $\alpha = \ordtype{B, R}$ आणि $\tuple{B, R} \in C$. म्हणून $\alpha
\in \alpha$, ही विसंगती मिळते.
\end{proof}

यावरून एक सखोल निष्कर्ष मिळतो:

\begin{thm}[\emph{$\ZF$ मध्ये}]
पुढील विधाने सममूल्य आहेत:
\begin{enumerate}
	\item\label{sth:choice:hartogs:equivwo} सुक्रमणाचे स्वयंसिद्धक
	\item\label{sth:choice:hartogs:equivcompare} कोणत्याही संच $A$ आणि $B$ साठी $\cardle{A}{B}$ किंवा
	$\cardle{B}{A}$
\end{enumerate}
\end{thm}

\begin{proof}
\emph{\ref{sth:choice:hartogs:equivwo} $\Rightarrow$ \ref{sth:choice:hartogs:equivcompare}.} $A$ आणि $B$ निश्चित करा. \ref{sth:choice:hartogs:equivwo} वापरल्यास $\tuple{A, R}$ आणि $\tuple{B, S}$ ही सुक्रमणे मिळतात. \ref{sth:ordinals:ordtype:thmOrdinalRepresentation} वापरून $f \colon
\alpha \to \tuple{A, R}$ आणि $g \colon \beta \to \tuple{B, S}$ ही समरूपणे असू द्या. \ref{sth:ordinals:basic:ordinalsaresubsets} नुसार $\alpha \subseteq \beta$ किंवा $\beta \subseteq \alpha$. $\alpha \subseteq \beta$ असेल, तर $\comp{f^{-1}}{g} \colon A \to B$ हे एकास-एक फलन आहे आणि म्हणून
$\cardle{A}{B}$; त्याचप्रमाणे $\beta \subseteq \alpha$ असेल, तर $\cardle{B}{A}$.

\emph{\ref{sth:choice:hartogs:equivcompare} $\Rightarrow$ \ref{sth:choice:hartogs:equivwo}.} $A$ निश्चित करा. \ref{sth:choice:hartogs:HartogsLemma} नुसार $\cardnless{\beta}{A}$ असलेली क्रमसूचक संख्या $\beta$ आहे. \ref{sth:choice:hartogs:equivcompare} वापरल्यास $\cardle{A}{\beta}$ मिळते. म्हणून एकास-एक फलन $f \colon A \to
\beta$ अस्तित्वात आहे, आणि $\Setabs{\tuple{a, b} \in A \times A}{f(a)
\in f(b)}$ हा क्रम परिभाषित करून त्या फलनाने $A$ च्या घटकांचे सुक्रमण करता येते.
\end{proof}
\noindent
याचे तात्काळ फलित असे: सुक्रमण अपयशी ठरल्यास काही संच त्यांच्या आकारमानाच्या दृष्टीने \emph{अक्षरशः अतुलनीय} असतात. म्हणून सुक्रमण अपयशी ठरल्यास अनंतपार संचांकांचे अंकगणित गुंतागुंतीचे होते. उदाहरणार्थ, $A$ आणि $B$ अनंत असल्यास $\cardeq{\cardeq{A \disjointsum B}{A \times
B}}{M}$, जेथे $M$ हा $A$ आणि $B$ यांपैकी मोठा संच आहे, ही कल्पना सोडावी लागेल (पाहा \ref{sth:card-arithmetic:simp:cardplustimesmax}). अडचण साधी आहे: $A$ आणि $B$ यांच्या आकारमानांची \emph{तुलनाच} करता येत नसेल, तर कोणता मोठा आहे असा प्रश्न विचारण्याला अर्थ नाही.







\section{सुक्रमणाचा प्रश्न}\label{sth:choice:woproblem:sec}

सुक्रमण स्वीकारतो की नाही, यावर बऱ्याच गोष्टी अवलंबून आहेत. पण या तत्त्वाची चर्चा मुख्यतः त्याच्याशी सममूल्य असलेल्या निवडीच्या स्वयंसिद्धकावर केंद्रित झाली आहे. म्हणून आता त्याकडे वळू (आणि ही सममूल्यता सिद्ध करू).

1883 मध्ये कँटरने सुक्रमणाच्या स्वयंसिद्धकाला पाठिंबा दर्शवला. त्याच्या मते तो “विचाराचा असा नियम आहे की जो मूलभूत, फलितांनी समृद्ध, आणि विशेषतः त्याच्या सार्वत्रिक युक्ततेमुळे उल्लेखनीय वाटतो” (पॉटर 2004, पृ.~243 मध्ये उद्धृत). पण शेवटी या “स्वयंसिद्धकाला” पुराव्याची गरज आहे, असे कँटरचे मत झाले. गणितज्ञ समुदायाचेही तसेच मत झाले.

1904 मध्ये झर्मेलोने या प्रश्नाचे “समाधान” केले. त्याचे समाधान स्पष्ट करण्यासाठी काही व्याख्या लागतील.

\begin{defn}
प्रत्येक $x \in \dom{f}$ साठी $f(x) \in x$ असेल तेव्हा आणि केवळ तेव्हाच फलन $f$ हे \emph{निवडफलन} असते. $f$ हे निवडफलन असून $\dom{f} = A \setminus \{\emptyset\}$ असेल तेव्हा आणि केवळ तेव्हाच $f$ हे \emph{$A$ साठीचे निवडफलन} आहे, असे म्हणतो.
\end{defn}

सहज अर्थाने, प्रत्येक अरिक्त संच $x \in A$ साठी, $A$ चे निवडफलन त्या $x$ मधून विशिष्ट घटक $f(x)$ \emph{निवडते}. मग निवडीचे स्वयंसिद्धक असे आहे:

\begin{axiom}[निवड]
	प्रत्येक संचाला निवडफलन असते.
\end{axiom}

निवडीवरून सुक्रमण मिळते आणि उलटही मिळते, हे झर्मेलोने दाखवले:

\begin{thm}[$\ZF$ मध्ये]\label{sth:choice:woproblem:thmwochoice}
सुक्रमण आणि निवड सममूल्य आहेत.
\end{thm}

\begin{proof}
\emph{डावीकडून उजवीकडे.} $A$ हा संचांचा संच असू द्या. युतीकरणाच्या स्वयंसिद्धकामुळे $\bigcup A$ अस्तित्वात असतो; म्हणून सुक्रमणामुळे $\bigcup A$ ला सुक्रमित करणारा $<$ आहे. आता अरिक्त सदस्यांसाठी $f(x) = \text{$x$ चा $<$-लघुतम घटक}$ असे घ्या. हे $A$ साठीचे निवडफलन आहे.

\emph{उजवीकडून डावीकडे.} $A$ निश्चित करा. रिक्त संचाचे सुक्रमण रिक्त क्रमाने मिळते; म्हणून पुढे आधारसंच अरिक्त मानू. निवडीमुळे $\Pow{A} \setminus \{\emptyset\}$ साठी निवडफलन $f$ आहे. अनंतपार पुनरावर्तन वापरून पुढील फलनाची रचना करू:
\begin{align*}
	g(0) &= f(A)\\
	g(\alpha) &=
		\begin{cases}
			\text{थांबा!{}} &\text{जर }A = \funimage{g}{\alpha}\\
			f(A \setminus \funimage{g}{\alpha}) & \text{अन्यथा}\\
		\end{cases}
\end{align*}
“थांबा!” ही सूचना अधिक लांबलचक व्याख्येचे संक्षिप्त रूप आहे. म्हणजे प्रथम $A =
\funimage{g}{\alpha}$ मिळाल्यावर प्रत्येक $\delta \geq \alpha$ साठी $g(\delta) = A$ घ्या. ही स्थिर पुढील मूल्ये आधीच्या निवडी बदलत नाहीत; पुराव्यातील एकास-एकपणाचा दावा थांबण्यापूर्वीच्या भागापुरता आहे. सुरुवातीला, याने केवळ एक \emph{पद} परिभाषित होते एवढीच खात्री आहे (\ref{sth:spine:recursion:transrecursionschema} नंतरच्या नोंदी पाहा); पण शेवटी थांबण्यापूर्वीचा भाग एका संच-प्रांतावर मर्यादित होतो आणि त्यामुळे आपल्याला फलन मिळते, हे दाखवू.

$f$ हे निवडफलन असल्यामुळे प्रत्येक $\alpha$ साठी (थांबण्यापूर्वी) $g(\alpha) =
f(A \setminus \funimage{g}{\alpha}) \in A \setminus
\funimage{g}{\alpha}$; म्हणजे $g(\alpha) \notin \funimage{g}{\alpha}$.
म्हणून $g(\alpha) = g(\beta)$ असल्यास $g(\beta) \notin
\funimage{g}{\alpha}$, म्हणजे $\beta \notin \alpha$; आणि त्याचप्रमाणे $\alpha \notin \beta$. म्हणून त्रिपर्यायाने $\alpha = \beta$. त्यामुळे या भागावर $g$ हे एकास-एक आहे.

आता पाहा की आपण थांबतो!{}; म्हणजे $A = g[\alpha]$ असलेली एक लघुतम क्रमसूचक संख्या $\alpha$ आहे. तसे नसल्यास $g$ हे एकास-एक असल्याने प्रत्येक क्रमसूचक संख्या $\alpha$ साठी $\cardle{\alpha}{A}$ मिळेल, जे \ref{sth:choice:hartogs:HartogsLemma} शी विसंगत आहे. त्यामुळे थांबण्यापूर्वीच्या भागाला मर्यादित केल्यावर $\ran{g} = A$ सुद्धा मिळते.

ही तथ्ये एकत्र केल्यावर, एखाद्या क्रमसूचक संख्येपासून $A$ पर्यंत $g$ हे एकास-एक आच्छादन आहे; येथे फक्त थांबण्यापूर्वीचा भाग घेतला आहे. आता $g$ वापरून $A$ चे सुक्रमण करता येते.
\end{proof}

म्हणून सुक्रमण आणि निवड हे एकत्रच स्वीकारले किंवा नाकारले जातात. पण प्रश्न उरतो: त्यांचा स्वीकार करावा की नकार?








\section{गणनीय निवड}\label{sth:choice:countablechoice:sec}

निवड किंवा सुक्रमण अजिबात न वापरता पुढील निष्कर्ष सहज सिद्ध करता येतो:

\begin{lem}[$\Zminus$ मध्ये]
प्रत्येक सांत संचाला निवडफलन असते.
\end{lem}

\begin{proof}
$a = \{b_1, \ldots, b_n\}$ असू द्या; सदस्यांची पुनरावृत्ती करू नका. सुलभतेसाठी प्रत्येक $b_i
\neq \emptyset$ आहे असे समजा. मग $c_1 \in b_1, \ldots, c_n \in b_n$ असलेल्या वस्तू $c_1, \ldots, c_n$ आहेत. आता \ref{sth:z:pairs:prop:pairsconsequences} नुसार $\{\langle b_1,
c_1\rangle , \ldots, \langle b_n, c_n\rangle\}$ हा संच अस्तित्वात आहे; आणि हे $a$ साठीचे निवडफलन आहे.
\end{proof}

पण अनंत संचांचा विचार केल्याबरोबर गोष्टी अधिक अस्पष्ट होतात. उदाहरणार्थ, वरील विधानाचा हा “किमान” विस्तार पाहा:

\begin{defish}
\emph{गणनीय निवड.} प्रत्येक \emph{गणनीय} संचाला निवडफलन असते.
\end{defish}

हे निवडीचे एक विशेषप्रकरण आहे. हे तत्त्व वापरत असल्याची स्पष्ट जाणीव नसतानाही ते बऱ्याचदा वापरले गेले होते, असे दिसते. पुढील दोन सुंदर उदाहरणे पाहा.\footnote{ही उदाहरणे पॉटर (2004, \S9.4) आणि लूका इन्कुर्वाती यांच्याकडून घेतली आहेत.}

\begin{ex}
एक स्वाभाविक विचार असा: कोणत्याही संच $A$ साठी $\cardle{\omega}{A}$ किंवा एखाद्या $n \in \omega$ साठी $\cardeq{A}{n}$. प्रत्येक संच सांत किंवा अनंत असतो, ही सहज कल्पना मांडण्याचा हा एक मार्ग आहे. कँटर आणि इतर अनेक गणितज्ञांनी पुराव्याशिवाय हा दावा केला. आपण सावध असल्यामुळे \ref{sth:cardinals:classing:generalinfinitycharacter} मध्ये तो सिद्ध केला. पण त्या पुराव्यात कोणत्याही संच $A$ चे सुक्रमण करता येते आणि म्हणून $\card{A}$ अस्तित्वात असल्याची खात्री आहे, असे मानून आपण $\ZFC$ मध्ये काम केले. म्हणजे आपण निवड स्पष्टपणे गृहीत धरली होती.

खरे तर डेडेकिंड (1888) यांनी या दाव्याचा स्वतःचा पुरावा पुढीलप्रमाणे दिला होता:

\begin{thm}[$\Zminus + \text{गणनीय निवड}$ मध्ये]
कोणत्याही $A$ साठी $\cardle{\omega}{A}$ किंवा एखाद्या $n \in \omega$ साठी $\cardeq{A}{n}$.
\end{thm}

\begin{proof}
प्रत्येक $n \in \omega$ साठी $\cardneq{A}{n}$ असे समजा. मग विशेषतः प्रत्येक $n < \omega$ साठी नेमके $\cardexpo{2}{n}$ घटक असलेला उपसंच $A_n \subseteq A$ आहे. $A_0, A_1, A_2,
\ldots$ हा अनुक्रम वापरून प्रत्येक $n$ साठी पुढील व्याख्या करू:
\[
	B_n = A_n \setminus \bigcup_{i < n} A_i.
\]
शून्यव्या निर्देशांकासाठी आधीचे युतीकरण रिक्त आहे आणि पहिला उपसंच अरिक्त आहे. धन निर्देशांकांसाठी पुढील मोजणी पाहा:
\begin{align*}
	\card{\bigcup_{i < n}A_i}
	&\leq \card{A_0} + \card{A_1} + \ldots + \card{A_{n-1}}\\
	&=1 + 2 + \ldots + 2^{n-1}\\
	& = 2^n - 1\\
	& < 2^n = \card{A_n}
\end{align*}
म्हणून प्रत्येक $B_n$ मध्ये किमान एक घटक $c_n$ आहे. शिवाय $B_n$ हे संच जोडीजोडीने वियुक्त आहेत; त्यामुळे $c_n = c_m$ असल्यास $n = m$. पण प्रत्येक $c_n
\in A$. म्हणून $f(n) = c_n$ हे फलन $\omega \to A$ असे एकास-एक फलन आहे.
\end{proof}
\noindent
गणनीय निवड वापरली असल्याची डेडेकिंडने नोंद केली नाही. पण \emph{तुम्हाला} तिचा वापर दिसला का? पुन्हा पाहा. (खरेच: \emph{पुन्हा पाहा}.)

या पुराव्यात गणनीय निवड दोनदा वापरली. $A_0$, $A_1$, $A_2$, \dots\@ या संचांचा अनुक्रम मिळवण्यासाठी ती एकदा वापरली. त्यानंतर प्रत्येक $B_n$ मधून $c_n$ हे घटक निवडण्यासाठी ती पुन्हा वापरली. शिवाय येथे कोणत्याही प्रकारची निवड न वापरता हा निष्कर्ष सिद्ध करता येत नाही. कोहेन (1966, पृ.~138) यांनी निवडीचे कोणतेही रूप नसल्यास हा निष्कर्ष अपयशी ठरतो, असे सिद्ध केले. म्हणजे $\omega$ शी \emph{अतुलनीय} असलेले संच आहेत, ही बाब $\ZF$ शी सुसंगत आहे.
\end{ex}

\begin{ex}
1878 मध्ये कँटरने म्हटले होते की गणनीय संचांचे गणनीय युतीकरण गणनीय असते. त्याने पुरावा दिला नाही; कदाचित तो पुरावा स्वयंस्पष्ट आहे, असे त्याला वाटले असावे. आपण सावध असल्यामुळे \ref{sth:card-arithmetic:simp:kappaunionkappasize} मध्ये या निष्कर्षाचे अधिक सामान्य रूप सिद्ध केले. पण आपल्या पुराव्यात निवड स्पष्टपणे गृहीत धरली होती. आणि कमी सामान्य निष्कर्षाचा खालील पुरावाही गणनीय निवड वापरतो.

\begin{thm}[$\Zminus + \text{गणनीय निवड}$ मध्ये]
प्रत्येक $n \in \omega$ साठी $A_n$ गणनीय असेल, तर $\bigcup_{n <
\omega} A_n$ गणनीय असतो.
\end{thm}

\begin{proof}
सामान्यतेला बाधा न आणता प्रत्येक $A_n \neq \emptyset$ आहे असे मानता येते. मग प्रत्येक $n \in \omega$ साठी आच्छादन $f_n \colon \omega
\to A_n$ आहे. $f(m, n) = f_n(m)$ घेऊन $f \colon \omega \times \omega \to \bigcup_{n <
\omega} A_n$ परिभाषित करा. $\omega
\times \omega$ गणनीय आहे (\ref{sfr:siz:zigzag:natsquaredenumerable}) आणि $f$ हे आच्छादन आहे, म्हणून निष्कर्ष मिळतो.
\end{proof}
\noindent
गणनीय निवडीचा वापर दिसला का? $f_0$, $f_1$, $f_2$, \dots\footnote{\ref{sth:card-arithmetic:simp:kappaunionkappasize} मध्ये “प्रत्येक $\beta \in \cardfont{a}$ साठी एकास-एक फलन $f_\beta$ निश्चित करा” अशी सूचना देताना निवडीचा असाच वापर झाला होता.} हा फलनांचा अनुक्रम निवडण्यासाठी ती वापरली आहे. पुन्हा, कोणतेही निवडतत्त्व नसल्यास हा निष्कर्ष अपयशी ठरू शकतो. विशेषतः फेफरमन आणि लेव्ही (1963) यांनी असे सिद्ध केले की $\ZF$ सुसंगत असेल, तर वास्तव संख्यांचा संच $\Real$ हा गणनीय संचांचे गणनीय युतीकरण आहे, हे विधानही $\ZF$ शी सुसंगत आहे. पण याच मुद्द्याचे रसेल यांनी केलेले अधिक मजेशीर वर्णन पाहा:
\begin{quote}
  एक कोट्यधीश जेव्हा बुटांची जोडी विकत घेत असे, तेव्हाच मोज्यांचीही जोडी घेत असे, इतर वेळी कधीच नाही. दोन्ही गोष्टी विकत घेण्याची त्याला इतकी हौस होती की शेवटी त्याच्याकडे बुटांच्या $\aleph_0$ जोड्या आणि मोज्यांच्या $\aleph_0$ जोड्या झाल्या\dots\@ बुटांत उजवा आणि डावा असा फरक करता येतो. त्यामुळे प्रत्येक जोडीतून एक बूट निवडता येतो: सर्व उजवे किंवा सर्व डावे बूट निवडू शकतो. पण मोज्यांच्या बाबतीत असे कोणतेही निवडतत्त्व सुचत नाही; आणि गुणक स्वयंसिद्धक [म्हणजे प्रत्यक्षात निवड] गृहीत धरल्याशिवाय, प्रत्येक जोडीतून एक मोजा असलेला वर्ग आहे याची खात्री देता येत नाही.
  (रसेल 1919, पृ.~126)
\end{quote}
म्हणजे मोज्यांच्या गणनीय संख्येने जोड्या असतील, तर मोजेही फक्त गणनीय संख्येने असतात, हे सिद्ध करण्यासाठी काही प्रकारची निवड लागते. गणनीय निवड अशा युतीकरणाची गणनीयता सिद्ध करण्यास पुरेशी आहे; पण निवडतत्त्वांचा अभाव असताना गणनीय संचांचे गणनीय युतीकरण गणनीय असण्यात अपयश येऊ शकते. येथे गणनीय निवडीशी उलटी सममूल्यता सिद्ध केलेली नाही.
\end{ex}

यातून धडा असा मिळतो की वापरणाऱ्यांना फारशी जाणीव नसतानाही गणनीय निवड वारंवार वापरली गेली. तात्त्विक प्रश्न असा आहे: गणनीय निवडीचे \emph{समर्थन} कसे करावे?

सहज समर्थनाचा एखादा प्रयत्न सांत वेळेत अनंत क्रियांचा क्रम पूर्ण करण्याच्या कल्पनेचा आधार घेऊ शकतो. पहिली निवड $\nicefrac{1}{2}$ मिनिटात, दुसरी $\nicefrac{1}{4}$ मिनिटात, \dots, आणि $n$ वी निवड $\nicefrac{1}{2^n}$ मिनिटात, \dots\@ केली असे समजा. मग $1$ मिनिटात निवडींचा $\omega$-अनुक्रम पूर्ण करून निवडफलन परिभाषित केले असेल.

पण असा विचारप्रयोग खरोखर काय सांगू शकतो? प्रथम, “निवड करणे” ही कल्पना अगदी शब्दशः घेतल्यावर तो अवलंबून आहे. दुसरे म्हणजे त्यात गणित सत्तामीमांसक शक्यतेशी बांधले जात असल्यासारखे दिसते.

अधिक महत्त्वाचे असे की यामुळे \emph{केवळ} गणनीय निवडीऐवजी निवडीच्या \emph{पूर्ण रूपाचे} समर्थन मिळणार नाही. कारण \emph{प्रत्येक} संचाला निवडफलन हवे असेल, तर “कोणत्याही क्रमसूचक लांबीचा अनंत क्रियाक्रम” पूर्ण करता आला पाहिजे. स्पष्ट सांगायचे तर ही कल्पना हास्यास्पद आहे.







\section{निवडीविषयीचे आंतरिक विचार}\label{sth:choice:justifications:sec}

मग व्यापक प्रश्न असा आहे की सुक्रमणाचे, निवडीचे किंवा सर्व संचांच्या आकारमानांच्या तुलनाक्षमतेचे---यांपैकी कोणतेही म्हणा---समर्थन करता येते का?

पुढे \emph{आंतरिक} समर्थनाचा एक प्रयत्न पाहू. \ref{sth:z:story:sec} मध्ये श्रेणीविषयी काही तत्त्वे मांडली होती. त्यांपैकी एक पुन्हा सांगण्याजोगे आहे:
\begin{enumerate}
	\item[] \stagesacc. कोणत्याही टप्पा $S$ साठी, आणि $S$ च्या \emph{आधी} तयार झालेल्या कोणत्याही संचांसाठी: नेमके ते संच घटक असलेला संच $S$ या टप्प्यावर तयार होतो. $S$ या टप्प्यावर दुसरे काहीही तयार होत नाही.
\end{enumerate}
खरे तर या तत्त्वाद्वारे किंवा यासारख्या तत्त्वाद्वारे निवडीच्या स्वयंसिद्धकाचे समर्थन करता येते, असे अनेक लेखकांनी सुचवले आहे. या पद्धतीचे थोडक्यात स्पष्टीकरण देऊ.

निवडीशी सममूल्य असलेले \emph{आणखी एक} तत्त्व देणाऱ्या साध्या निष्कर्षापासून सुरुवात करू:

\begin{thm}[$\ZF$ मध्ये]\label{sth:choice:justifications:choiceset}
निवड पुढील तत्त्वाशी सममूल्य आहे. $A$ चे घटक वियुक्त आणि अरिक्त असतील, तर प्रत्येक $x \in A$ साठी $C
\cap x$ एकघटकी असलेला $C$ अस्तित्वात आहे. (अशा $C$ ला $A$ साठीचा {निवडसंच} म्हणतो.)
\end{thm}

या निष्कर्षाचा पुरावा सरळ आहे आणि वाचकासाठी सराव म्हणून सोडला आहे.

\begin{prob}
\ref{sth:choice:justifications:choiceset} सिद्ध करा. अडचण आल्यास (पॉटर 2004, पृ.~242--3) मध्ये पुरावा पाहता येईल.
\end{prob}

मुख्य मुद्दा असा की $A$ साठीच्या निवडसंचातून आधारकुटुंबाच्या युतीकरणाबाहेरील घटक वगळल्यावर तो $A$ साठीच्या निवडफलनाची प्रतिमा ठरतो. म्हणून निवडीचे समर्थन करण्यासाठी निवडसंचांच्या अस्तित्वाच्या भाषेतील तिच्या सममूल्य मांडणीचे समर्थन करण्याचा प्रयत्न करता येईल. आता नेमके तेच करू.

$A$ चे घटक वियुक्त आणि अरिक्त असू द्या. \stageshier{} नुसार (पाहा \ref{sth:z:story:sec}) $A$ कोणत्या तरी टप्पा $S$ वर तयार होतो. $\bigcup A$ चे सर्व घटक $S$ टप्प्याच्या आधी उपलब्ध असतात. आता \stagesacc{} नुसार, $S$ च्या आधी तयार झालेल्या \emph{कोणत्याही} संचांचे नेमके तेच घटक असलेला संच तयार होतो. दुसऱ्या शब्दांत: आधी उपलब्ध झालेल्या संचांचे प्रत्येक \emph{शक्य} संकलन $S$ वर अस्तित्वात असते. पण $A$ साठीच्या निवडसंचात एकत्र करता येतील अशा वस्तू निवडणे नक्कीच \emph{शक्य} आहे, असा या समर्थनाचा दावा आहे; ते $\bigcup A$ च्या अगदी विशिष्ट उपसंचातच येतील. म्हणून आवश्यक असा एखादा निवडसंच अस्तित्वात आहे, असे या युक्तिवादातून म्हटले जाते.

आंतरिक आधारावर निवडीचे समर्थन करण्याचा हा \emph{अतिशय} थोडक्यात मांडलेला प्रयत्न आहे. ही कल्पना पुढे नेण्यासाठी पॉटर यांनी (2004, \S14.8) केलेली सुबक मांडणी वाचावी.







\section{बानाख--टार्स्की विरोधाभास}\label{sth:choice:banach:sec}

बूलोस यांनी प्रतिस्थापनाचे समर्थन करण्याचा प्रयत्न केला तसा, \emph{बाह्य} विचारांचा आधार घेऊन निवडीचे समर्थन करण्याचाही प्रयत्न करता येईल (पाहा \ref{sth:replacement:extrinsic:sec}). अखेर निवड स्वीकारल्याने अनेक इष्ट परिणाम मिळतात: प्रत्येक संचांकाची तुलना करता येते, प्रत्येक संचाचे सुक्रमण करता येते, संचांकांना विशिष्ट प्रकारच्या क्रमसूचक संख्या मानता येते, इत्यादी.

पण कधी कधी निवडीचे \emph{अनिष्ट} परिणाम आहेत, असा दावा केला जातो. मुख्यतः (बानाख आणि टार्स्की 1924) यांच्या एका निष्कर्षामुळे हा दावा केला जातो.

\begin{thm}[बानाख--टार्स्की विरोधाभास ($\ZFC$ मध्ये)]
त्रिमितीय अवकाशातील धन त्रिज्येच्या कोणत्याही घनगोलाचे सांत संख्येने तुकडे करता येतात; घूर्णन आणि स्थानांतरणाने त्यांची पुनर्जुळणी करून त्या घनगोलाच्या दोन प्रती तयार करता येतात.
\end{thm}
\noindent
पहिल्या दृष्टीने हे थोडे थक्क करणारे आहे. या दोन घनगोलांचे घनफळ मूळ घनगोलाच्या \emph{दुप्पट} आहे. पण दृढ गती---घूर्णन आणि स्थानांतरण---घनफळ बदलत नाहीत. म्हणून बानाख--टार्स्कीमुळे जादूने नवे द्रव्य अस्तित्वात आणता येते, असे दिसते.

यापुढे आणखी आश्चर्य आहे.\footnote{पाहा वॅगन (2016, प्रमेय 3.12).} यासारखा युक्तिवाद दाखवतो की वाटाण्याचे सांत संख्येने तुकडे करून घूर्णन आणि स्थानांतरणाने त्यांची पुनर्जुळणी केल्यास बिग बेनच्या आकाराची आणि आकारमानाची वस्तू तयार करता येते.

पण केवळ $\ZF$ मध्ये यापैकी काहीही सिद्ध होत नाही.\footnote{तरी निवडीपेक्षा काटेकोरपणे दुर्बल तत्त्वांनी बानाख--टार्स्की सिद्ध करता येतो; पाहा वॅगन (2016, पृ.~303).} म्हणून निवड नाकारावी की या “विरोधाभासासह” राहायला शिकावे, हा निर्णय घ्यावा लागतो.

या “विरोधाभासासह” राहायला शिकावे, असे आपण सुचवू. खरे तर हा फारसा विरोधाभास आहे, असे आम्हाला वाटत नाही. विशेषतः पुढील निष्कर्षांपेक्षा तो कमी किंवा अधिक विरोधाभासी का आहे, हे दिसत नाही:\footnote{पॉटर (2004, पृ.~276--7), वेस्टन (2003, पृ.~16) आणि वॅगन (2016, 31, 308--9) इतर उदाहरणांनी असेच मुद्दे मांडतात.}
\begin{enumerate}
	\item $(0,1)$ या अंतरालात जितके बिंदू आहेत, तितकेच $\Real$ मध्ये आहेत.
	\\\emph{पुरावा}: $\tan(\pi(r-\nicefrac{1}{2}))$ पाहा.
	\item रेषेत जितके बिंदू आहेत, तितकेच चौरसात आहेत.
	\\पाहा \ref{his:set:pathology:sec} आणि \ref{his:set:cantorplane:sec}.
	\item अवकाशभरण करणाऱ्या वक्ररेषा आहेत.
	\\पाहा \ref{his:set:pathology:sec} आणि \ref{his:set:hilbertcurve:sec}.
\end{enumerate}
या तिन्ही निष्कर्षांना निवडीची गरज नाही. आज आपण त्यांना गणिताचे आश्चर्यकारक, सुंदर भाग मानतो. बानाख--टार्स्की विरोधाभासाविषयीही तीच भूमिका घ्यावी, असे वाटू शकते.

येथे एक तांत्रिक निरीक्षण आवश्यक आहे; पण त्यासाठी शांतपणे विचार करणे पुरेसे आहे. दृढ गती घनफळ जतन करतात. त्यामुळे घनगोलाचे जे पाच\footnote{विरोधाभासाचे विधान “सांत संख्येने तुकडे” असे केले होते. खरे तर रॉबिन्सन (1947) यांनी \emph{पाच} तुकड्यांत (त्याहून कमी तुकड्यांत नाही) हे विघटन करता येते, असे सिद्ध केले. पुराव्यासाठी वॅगन (2016, पृ.~66--7) पाहा.} तुकडे केले जातात, ते सर्व \emph{मापनीय} असू शकत नाहीत. साधारणपणे, या स्वतंत्र तुकड्यांना घनफळ देण्यात अर्थ नाही. त्यांना चित्रात दाखवता न येणारे बिंदूंचे “अनंत विखुरणे” समजावे. अशा “अनंत विखुरलेल्या” संचांची कल्पना करणे विचित्र वाटू शकते. पण आधी उपलब्ध झालेल्या संचांची \emph{सर्व शक्य} संकलने तयार करावीत, या \stagesacc{} मध्ये व्यक्त झालेल्या निर्देशावरून त्यांचे अस्तित्व मिळत असल्यासारखे दिसते.

हे पटत नसेल, तर बानाख--टार्स्की विरोधाभास स्वीकारण्याच्या बाजूने शेवटचा एक बाह्य युक्तिवाद पाहा. त्यावरून गणितातील उत्कृष्ट विनोद लगेच मिळतो:
\begin{enumerate}
	\item[] \emph{प्रश्न}. “बानाख--टार्स्की”च्या अक्षरांची पुनर्रचना केल्यावर काय मिळते?
	\item[] \emph{उत्तर}. “बानाख--टार्स्की बानाख--टार्स्की”.
\end{enumerate}







\section{परिशिष्ट: वितालीचा विरोधाभास}\label{sth:choice:vitali:sec}

बानाख--टार्स्कीची रचना स्वीकारार्ह आहे की नाही, हे नीट समजून घेण्यासाठी तिचा \emph{पुरावा} तपासला पाहिजे. दुर्दैवाने त्यासाठी येथे मांडता येईल त्याहून बरेच अधिक बीजगणित लागेल. तरी पुढील निष्कर्षावर लक्ष केंद्रित करून बानाख--टार्स्कीच्या पुराव्यावर प्रकाश टाकतील अशा काही थोडक्यात नोंदी देता येतील.\footnote{अधिक पूर्ण विवेचनासाठी (वेस्टन 2003) किंवा (वॅगन 2016) पाहा.}

\begin{thm}[वितालीचा विरोधाभास ($\ZFC$ मध्ये)]\label{sth:choice:vitali:vitaliparadox}
धन त्रिज्येच्या कोणत्याही वर्तुळाचे गणनीय संख्येने तुकडे करता येतात; घूर्णन आणि स्थानांतरणाने त्यांची पुनर्जुळणी करून त्या वर्तुळाच्या दोन प्रती तयार करता येतात.
\end{thm}

वितालीचा विरोधाभास बानाख--टार्स्की विरोधाभासापेक्षा सिद्ध करणे खूप सोपे आहे. वितालीने 1905 मध्ये केलेल्या अमापनीय संचाच्या रचनेवरून तो मिळत असल्यामुळे त्याला “वितालीचा विरोधाभास” म्हटले आहे. पण विताली आणि बानाख--टार्स्की यांच्या पुराव्यांचे संचसैद्धान्तिक पैलू अगदी सारखे आहेत. मुख्य फरक एवढाच की बानाख--टार्स्कीमध्ये \emph{सांत} विघटन आहे, तर वितालीमध्ये \emph{गणनीय अनंत} विघटन आहे. वेस्टन (2003) यांच्या शब्दांत, वितालीचा विरोधाभास “बानाख--टार्स्की विरोधाभासाइतका नक्कीच थक्क करणारा नाही; पण अगदी ‘साध्या’ परिस्थितीतही भूमितीय विरोधाभास येऊ शकतात, हे तो दाखवतो.”

वितालीचा विरोधाभास वर्तुळ या द्विमितीय आकृतीविषयी आहे. म्हणून $\Real^2$ या प्रतलावर काम करू. मूलबिंदूभोवती पूर्ण फेरीच्या \emph{परिमेय} भागांनुसार, $[0,2\pi)$ या मर्यादेत रेडियनमूल्य असलेली दक्षिणावर्त घूर्णने $\rotationsgroup$ मध्ये घ्या. येथे पूर्ण फेरीचा परिमेय भाग अभिप्रेत आहे; परिमेय संख्येइतके रेडियन नव्हेत. $\rotationsgroup$ विषयी पुढील बीजगणितीय तथ्ये पाहा (विधान समजले नाही, तर पुराव्यात त्याचा अर्थ स्पष्ट होईल):

\begin{lem}\label{sth:choice:vitali:rotationsgroupabelian}
फलनांच्या संयोजनाखाली $\rotationsgroup$ हा अबेली, म्हणजे क्रमनिरपेक्ष, {गट} होतो.
\end{lem}

\begin{proof}
$0$ रेडियनचे घूर्णन $0_{\rotationsgroup}$ असे लिहिल्यास हा $\rotationsgroup$ चा अविकारी घटक आहे. कारण प्रत्येक $\rho \in \rotationsgroup$ साठी $\comp{0_{\rotationsgroup}}{\rho} = \comp{\rho}{0_{\rotationsgroup}} =
\rho$.

प्रत्येक घटकाला व्यस्त असतो. शून्य घूर्णनाचा व्यस्त तोच असतो. इतर प्रकरणांत, $\rho \in \rotationsgroup$ हे $r$ रेडियनचे घूर्णन असेल, तर $\rho^{-1} \in \rotationsgroup$ हे $2\pi - r$ रेडियनचे घूर्णन असते. त्यामुळे $\rho \circ \rho^{-1} = 0_\rotationsgroup$.

संयोजनात सहचारिता आहे: कोणत्याही $\rho, \sigma, \tau \in
\rotationsgroup$ साठी $\comp{\rho}{(\comp{\sigma}{\tau})} =
\comp{(\comp{\rho}{\sigma})}{\tau}$.

संयोजन क्रमनिरपेक्ष आहे: कोणत्याही $\rho, \sigma \in \rotationsgroup$ साठी $\comp{\rho}{\sigma} =
\comp{\sigma}{\rho}$.
\end{proof}

खरे तर $\rotationsgroup$ या गटाचे दोन भाग करून प्रत्येक भागावरून पूर्ण गट पुन्हा मिळवता येतो:

\begin{lem}\label{sth:choice:vitali:disjointgroup}
$\rotationsgroup$ चे $\rotationsgroup_{1}$ आणि $\rotationsgroup_{2}$ अशा दोन वियुक्त संचांत विभाजन करता येते; दोन्ही $\rotationsgroup$ चे जनक संच आहेत.
\end{lem}

\begin{proof}
पूर्ण फेरीच्या परिमेय भागांशी जुळणारी आणि रेडियनमूल्य $[0, \pi)$ मध्ये असलेली घूर्णने $\rotationsgroup_{1}$ मध्ये घ्या; आणि $\rotationsgroup_{2} = \rotationsgroup
\setminus \rotationsgroup_1$ घ्या. प्राथमिक बीजगणिताने
$\Setabs{\comp{\rho}{\rho}}{\rho \in \rotationsgroup_1} =
\rotationsgroup$. $\rotationsgroup_2$ साठीही तसाच निष्कर्ष मिळतो.
\end{proof}

गटांविषयीचे हे तथ्य \ref{sth:choice:vitali:vitaliparadox} सिद्ध करण्यासाठी वापरू. $\onesphere$ हे एकक वर्तुळ, म्हणजे प्रतलाच्या मूलबिंदूपासून नेमके $1$ एकक अंतरावरील बिंदूंचा संच, म्हणजे $\Setabs{\tuple{r,s} \in \Real^2}{\sqrt{r^2+s^2}=1}$ असू द्या. $\onesphere$ वरील पुढील संबंध पाहून $\onesphere$ चे भाग करू:
\[
	r \sim s \emph{ असेल तेव्हा आणि केवळ तेव्हाच }(\exists \rho \in \rotationsgroup)\rho(r) = s.
\]
म्हणजे मूलबिंदूभोवती पूर्ण फेरीच्या परिमेय भागाचे घूर्णन करून एका बिंदूपासून दुसऱ्याकडे जाता येत असेल तेव्हा आणि केवळ तेव्हाच $\onesphere$ चे ते बिंदू या संबंधाने जोडलेले आहेत. अपेक्षेप्रमाणे:

\begin{lem}
$\sim$ हा तुल्यता संबंध आहे.
\end{lem}

\begin{proof}
\ref{sth:choice:vitali:rotationsgroupabelian} वापरून हे तात्काळ मिळते.
\end{proof}

आता निवड वापरून $\sim$ नुसार $\onesphere$ च्या प्रत्येक तुल्यतावर्गातून नेमका एक घटक असलेला संच $C$ मिळवू. म्हणजे तुल्यतावर्गांच्या पुढील संचावर निवडफलन $f$ घेऊ:\footnote{$\rotationsgroup$ हे गणनीय असल्यामुळे $E$ चा प्रत्येक घटक गणनीय आहे. $\onesphere$ हे अगणनीय असल्यामुळे \ref{sth:choice:vitali:vitalicover} आणि \ref{sth:card-arithmetic:simp:kappaunionkappasize} नुसार $E$ हे अगणनीय आहे. म्हणून येथे \emph{अ}गणनीय निवड वापरली आहे.}
\[
	E = \Setabs{\equivrep{r}{\sim}}{r \in \onesphere},
\]
आणि $C = \ran{f}$ घ्या. प्रत्येक घूर्णन $\rho \in \rotationsgroup$ साठी, $C$ च्या प्रत्येक बिंदूवर $\rho$ लागू करून मिळणाऱ्या बिंदूंचा संच $\funimage{\rho}{C}$ आहे. पुढील दोन निष्कर्ष दाखवतात की हे संच वर्तुळ पूर्णपणे व्यापतात आणि एकमेकांवर आच्छादित होत नाहीत:

\begin{lem}\label{sth:choice:vitali:vitalicover}
$\onesphere = \bigcup_{\rho \in \rotationsgroup} \funimage{\rho}{C}$.
\end{lem}

\begin{proof}
$s \in \onesphere$ निश्चित करा. $r \in \equivrep{s}{\sim}$ असलेला $r \in C$ आहे; म्हणजे $r \sim s$, म्हणजे एखाद्या $\rho \in \rotationsgroup$ साठी $\rho(r) = s$.
\end{proof}

\begin{lem}\label{sth:choice:vitali:vitalinooverlap}
$\rho_1 \neq \rho_2$ असल्यास $\funimage{\rho_1}{C} \cap \funimage{\rho_2}{C} = \emptyset$.
\end{lem}

\begin{proof}
$s \in \funimage{\rho_{1}}{C} \cap \funimage{\rho_{2}}{C}$ असे समजा. मग एखाद्या $r_{1}, r_{2} \in C$ साठी $s = \rho_{1}(r_{1}) = \rho_{2}(r_{2})$.
त्यामुळे $\rho^{-1}_2(\rho_1(r_1)) = r_2$, आणि $\comp{\rho_1}{\rho^{-1}_2} \in \rotationsgroup$; म्हणून $r_{1} \sim
r_{2}$. $\sim$ नुसार प्रत्येक तुल्यतावर्गातून $C$ नेमका एक घटक निवडतो, म्हणून $r_1 = r_2$. त्यामुळे $s = \rho_1(r_1) = \rho_2(r_1)$ आणि म्हणून $\rho_1 = \rho_2$.
\end{proof}

आता आधीची बीजगणितीय तथ्ये वर्तुळाला लागू करू:

\begin{lem}\label{sth:choice:vitali:pseudobanachtarski}
$\onesphere$ चे $D_{1}$ आणि $D_{2}$ अशा दोन वियुक्त संचांत विभाजन करता येते, ज्यामध्ये $D_{1}$ चे गणनीय संख्येने भाग करून ते घूर्णित केल्यास $\onesphere$ ची प्रत तयार होते (आणि $D_{2}$ साठीही तसेच).
\end{lem}

\begin{proof}
\ref{sth:choice:vitali:disjointgroup} मधील $\rotationsgroup_{1}$ आणि $\rotationsgroup_{2}$ वापरून पुढील व्याख्या करा:
\begin{align*}
	D_{1} &= \bigcup_{\rho \in \rotationsgroup_1} \funimage{\rho}{C} &
	D_{2} &= \bigcup_{\rho \in \rotationsgroup_2} \funimage{\rho}{C}
\end{align*}
\ref{sth:choice:vitali:vitalicover} नुसार हे $\onesphere$ चे विभाजन आहे, आणि \ref{sth:choice:vitali:vitalinooverlap} नुसार $D_1$ आणि $D_2$ वियुक्त आहेत. रचनेनुसार $D_1$ चे गणनीय संख्येने भाग करता येतात: प्रत्येक $\rho \in \rotationsgroup_1$ साठी $\funimage{\rho}{C}$ हा एक भाग. \ref{sth:choice:vitali:disjointgroup} आणि \ref{sth:choice:vitali:vitalicover} नुसार $\onesphere = \bigcup_{\rho \in
\rotationsgroup}\funimage{\rho}{C} = \bigcup_{\rho \in
\rotationsgroup_1}\funimage{(\comp{\rho}{\rho})}{C}$ असल्यामुळे हे भाग घूर्णित करून $\onesphere$ ची प्रत तयार करता येते. $D_2$ साठीही तोच युक्तिवाद लागू होतो. \end{proof}\noindent यावरून वितालीचा विरोधाभास लगेच मिळतो. वर्तुळाचे गणनीय संख्येने तुकडे (विविध $\funimage{\rho}{C}$) करून दृढ गतीने ते हलवल्यास $\onesphere$ पासून $\onesphere$ च्या \emph{दोन} प्रती तयार करता येतात. फक्त $D_1$ चा प्रत्येक भाग घूर्णित करा; आणि $D_2$ च्या प्रत्येक भागाचे आधी दुसरीकडे स्थानांतरण करून नव्या केंद्राभोवती घूर्णन करा.

पुराव्याची पद्धत पुन्हा पाहू. प्रतलावरील घूर्णन गटाच्या काही बीजगणितीय तथ्यांपासून सुरुवात केली. हा गट वापरून $\onesphere$ चे तुल्यतावर्गांत विभाजन केले. त्यानंतर प्रत्येक वर्गातून घटक निवडण्यासाठी निवड वापरून “विरोधाभास” मिळवला.

बानाख--टार्स्की सिद्ध करण्यासाठी नेमकी हीच पद्धत वापरतो. मुख्य फरक असा की त्यासाठीची बीजगणितीय तथ्ये वितालीच्या पुराव्यातील तथ्यांपेक्षा बरीच गुंतागुंतीची आहेत. पण त्या बीजगणितीय तथ्यांचा निवडीशी संबंध नाही. त्यांचा थोडक्यात सारांश देऊ.

बानाख--टार्स्की सिद्ध करण्यासाठी आधी \ref{sth:choice:vitali:disjointgroup} सारखा निष्कर्ष मिळवतो: दोन जनकांवरील \emph{मुक्त गटाचे} चार भाग करता येतात, आणि सहज अर्थाने ते “हलवून” पूर्ण गटाच्या दोन प्रती मिळवता येतात.\footnote{\emph{चार} भाग वापरता येतात, हा निष्कर्ष (रॉबिन्सन 1947) यांचा आहे. अलीकडील पुराव्यासाठी वॅगन (2016, प्रमेय 5.2) पाहा. गटाचे भाग “हलवणे” असे म्हणताना वेस्टन (2003, पृ.~3) यांच्या वर्णनाचा आधार घेतला आहे.} त्यानंतर $\Real^3$ च्या मूलबिंदूभोवतीची दोन विशिष्ट घूर्णने वापरून घूर्णनांचा मुक्त गट $F$ तयार करता येतो, हे दाखवतो.\footnote{पाहा वॅगन (2016, प्रमेय 2.1).} (अजून निवड वापरलेली नाही.) आता $F$ मधील घूर्णनाने एका बिंदूपासून दुसरा मिळत असेल तेव्हा आणि केवळ तेव्हाच गोलकाच्या पृष्ठभागावरील ते बिंदू “सदृश” मानतो. मग “सदृश” बिंदूंच्या प्रत्येक तुल्यतावर्गातून नेमका एक बिंदू निवडण्यासाठी \emph{निवड वापरतो}. \ref{sth:choice:vitali:pseudobanachtarski} प्रमाणे $F$ चे विभाजन पृष्ठभागाला लागू करण्याची कल्पना अशी की चार भाग “हलवून” पृष्ठभागाच्या दोन प्रती मिळवाव्यात. मात्र या त्रिमितीय रचनेत अक्षांवरील काही बिंदू घूर्णनांनी स्थिर राहतात; अशा बिंदूंचा स्वतंत्र विचार आणि आवश्यक अतिरिक्त युक्तिवाद या संक्षिप्त रूपरेषेत दिलेला नाही. म्हणून ही पूर्ण पुराव्याची जागा घेऊ शकत नाही. संबंधित निष्कर्ष (हाउसडॉर्फ 1914) असा मांडला जातो:

\begin{thm}[हाउसडॉर्फचा विरोधाभास ($\ZFC$ मध्ये)]
धन त्रिज्येच्या कोणत्याही गोलकाच्या पृष्ठभागाचे सांत संख्येने तुकडे करता येतात; घूर्णन आणि स्थानांतरणाने त्यांची पुनर्जुळणी करून त्या पृष्ठभागाच्या दोन वियुक्त प्रती तयार करता येतात.
\end{thm}

पूर्ण बानाख--टार्स्की प्रमेय मिळवण्यासाठी (जे फक्त गोलकाच्या पृष्ठभागाविषयी नाही, तर अंतर्भागाविषयीही आहे) आणखी काही बीजगणितीय युक्त्या लागतात. तरी लेखकाच्या मते मुख्य बीजगणितीय काम आधीच्या टप्प्यावर आहे. म्हणून वेस्टन लिहितात:
\begin{quote}
  [\ldots] मुक्त गटांविषयीचा निष्कर्ष बानाख--टार्स्की विरोधाभासाच्या पुराव्यातील \emph{मुख्य टप्पा} आहे. या दृष्टीने बानाख--टार्स्की विरोधाभास हा $\Real^3$ विषयीच्या विधानापेक्षा [$\Real^3$ मधील स्थानांतरणे आणि घूर्णने यांच्या] गटाच्या गुंतागुंतीविषयीचे विधान आहे. (वेस्टन 2003, पृ.~16)
\end{quote}
म्हणजे बानाख--टार्स्कीप्रमाणे \emph{सांत} विघटन मिळते की वितालीप्रमाणे \emph{गणनीय अनंत} विघटन मिळते, हे दोन किंवा तीन मितींत काम करताना गटांविषयी लागू होणाऱ्या काही तथ्यांवर ठरते.

हे शेवटचे निरीक्षण \ref{sth:choice:banach:sec} च्या अखेरचा विनोद थोडा बिघडवते. लेखन द्विमितीय असल्यामुळे “बानाख--टार्स्की बानाख--टार्स्की” अशी पुनर्जुळणी करायची असेल, तर “बानाख--टार्स्की”चे गणनीय \emph{अनंत} तुकडे करावे लागतील. विनोद दुरुस्त करण्यासाठी त्रिमितीय लेखन करावे. हा सराव वाचकासाठी सोडला आहे.

शेवटची एक नोंद. \ref{sth:choice:banach:sec} मध्ये गोलकाचे मिळणारे “तुकडे” सर्व \emph{मापनीय} असू शकत नाहीत; त्यात चित्रात दाखवता न येणारे “अनंत विखुरणे” असते, असे म्हटले होते. \ref{sth:choice:vitali:pseudobanachtarski} मिळवण्यासाठी केलेल्या निवडीच्या वापराविषयीही हेच लागू होते. हे स्पष्ट करणे उपयुक्त आहे.

पुन्हा थोडी पार्श्वभूमी सांगावी लागेल (ही \emph{केवळ} रूपरेषा आहे; \emph{मापाविषयी} पाठ्यपुस्तकातील विवेचन पाहावे). येथे वापरल्या जाणाऱ्या सांत मापात, संच $X$ वरील एखाद्या “$\sigma$-बीजगणितातील” प्रत्येक $E$ ला ऋणेतर मूल्य $\mu(E) \in
\Real$ दिले जाते; हे $X$ चे माप ठरवते. तपशील येथे आवश्यक नाहीत; पण फलन $\mu$ ने गणनीय योगशीलतेचे तत्त्व पाळले पाहिजे: वियुक्त संचांच्या गणनीय युतीकरणाचे माप त्यांच्या स्वतंत्र मापांची बेरीज असते. म्हणजे $X_n$ हे संच वियुक्त असताना $\mu(\bigcup_{n < \omega} X_n) = \sum_{n < \omega}\mu(X_n)$.
संच “अमापनीय” आहे याचा अर्थ तो या निश्चित मापाच्या मापनीय क्षेत्रात येत नाही; कोणतेही माप कधीच देता येत नाही, असा सार्वत्रिक अर्थ नाही. आता आधीचा $\rotationsgroup$ वापरू:

\begin{cor}[विताली]
$\mu$ हे असे माप असू द्या की $\mu(\onesphere) = 1$, आणि $X$ व $Y$ एकरूप असतील, तर $\mu(X) = \mu(Y)$. मग प्रत्येक $\rho \in
\rotationsgroup$ साठी $\funimage{\rho}{C}$ अमापनीय आहे.
\end{cor}

\begin{proof}
विरोधसिद्धीसाठी तसे नाही असे समजा. मग एखाद्या $\sigma \in \rotationsgroup$ आणि $r \in \Real$ साठी $\mu(\funimage{\sigma}{C}) =
r$ असू द्या. कोणत्याही $\rho \in \rotationsgroup$ साठी, $\funimage{\rho}{C}$ आणि $\funimage{\sigma}{C}$ एकरूप आहेत; म्हणून प्रत्येक $\rho \in \rotationsgroup$ साठी $\mu(\funimage{\rho}{C}) = r$.
\ref{sth:choice:vitali:vitalicover} आणि \ref{sth:choice:vitali:vitalinooverlap} नुसार $\onesphere =
\bigcup_{\rho \in \rotationsgroup}\funimage{\rho}{C}$ हे जोडीजोडीने वियुक्त संचांचे गणनीय युतीकरण आहे. म्हणून गणनीय योगशीलतेनुसार $\mu(\onesphere) = 1$ ही प्रत्येक $\funimage{\rho}{C}$ च्या मापांची बेरीज आहे; म्हणजे
\[
	1 = \mu(\onesphere) = \sum_{\rho \in \rotationsgroup}\mu(\funimage{\rho}{C}) = \sum_{\rho \in \rotationsgroup}r
\]
पण $r = 0$ असल्यास $\sum_{\rho \in \rotationsgroup}r = 0$, आणि $r
> 0$ असल्यास $\sum_{\rho \in \rotationsgroup}r = \infty$.
\end{proof}



\part{पद्धती}\label{mth:part}
\begin{editorial}
  या भागात सर्वसाधारण आणि पद्धतिविषयक साहित्य आहे. विशेषतः गणिताचा विद्यार्थी नसलेल्या वाचकाला अपरिचित असू शकणाऱ्या विविध सिद्धता-पद्धतींची स्पष्टीकरणे आहेत. या स्रोतभागात सध्या सिद्धता कशा लिहाव्यात यावरील एक प्रकरण आणि विगमनावरील एक प्रकरण आहे. त्यांत आणखी विभाग व सराव जोडणे, तसेच गणितीय संज्ञांविषयी एक प्रकरण तयार करणे नियोजित आहे.
\end{editorial}
